\section{Conclusões}\label{sec:conclusions}

A investigação foi formulada para quantificar o efeito da representação por classes sobre a previsão de deslocamentos e o consumo de madeira, utilizando a base experimental de \textcite{Wolenski2020}. O protocolo distingue a conversão de resistência média em característica, a substituição da rigidez e o redimensionamento. A conclusão sobre atendimento ao estado limite de serviço será baseada na reavaliação do projeto por classe, e não apenas no sinal da diferença de rigidez ou de volume.

\begin{fillblock}
\noindent\textbf{[CONCLUSÕES QUANTITATIVAS PENDENTES]} Após concluir as análises, responder com valores e limites de validade:
\begin{enumerate}[nosep]
\item Quanto a conversão simplificada altera o enquadramento em relação aos valores característicos publicados?
\item Qual o erro de previsão de deslocamento introduzido pela representação da rigidez por classe, sob geometria fixa?
\item Quanto o redimensionamento altera o volume de madeira e em quais casos o projeto por classe excede o limite de serviço quando reavaliado com a rigidez experimental?
\item Como essas diferenças dependem do vão, das restrições ativas e das hipóteses de propriedades complementares?
\item Em quais condições os resultados justificam caracterização adicional da rigidez para a aplicação estudada?
\end{enumerate}
Não concluir previamente pela inadequação das classes, pela existência de um vão crítico ou pela superioridade de um sistema normativo. Efeitos pequenos, ausência de excedências e ausência de migração do estado limite governante também são desfechos admissíveis.
\end{fillblock}
